\chapter{Analyse et spécification des besoins}

\section*{Introduction}
Ce chapitre formalise l'analyse métier du système PMS. Il identifie les acteurs, recense les
besoins fonctionnels et non fonctionnels, présente les règles de gestion ainsi que le modèle
d'autorisation dynamique qui structure toute la plateforme, et détaille le modèle financier et les
indicateurs de performance (KPI) extraits de la réalité métier. Cette spécification constitue le
socle des phases de conception ultérieures.

\section{Identification des acteurs}
Le système repose sur quatre rôles principaux, complétés par des acteurs secondaires (le système
lui-même et l'outil externe de suivi des temps KIMAI).

\begin{longtable}{|m{3cm}|m{6.5cm}|m{5cm}|}
\hline
\textbf{Acteur} & \textbf{Mission} & \textbf{Limite (ne peut pas)} \\
\hline
\endhead
\endfoot
\endlastfoot
\hline
Administrateur & Administration technique et fonctionnelle~: utilisateurs, rôles, permissions, TCC,
paramétrage. & Piloter les projets~; accéder aux données financières des projets. \\
\hline
Directeur & Gouvernance stratégique~: création des projets, budgets, affectation des chefs de
projet, supervision du portefeuille, KPI exécutifs. & Gérer les opérations quotidiennes (délégué,
non interdit). \\
\hline
Chef de Projet & Utilisateur opérationnel principal~: équipes, plan de charge, facturation,
missions, suivi des coûts et de la rentabilité. & Gérer les utilisateurs, les rôles, le TCC et le
paramétrage. \\
\hline
Développeur & Saisie des charges réelles~; consultation de ses affectations et missions. & Accéder
à toute donnée financière (budget, marge, EAC, facturation). \\
\hline
\caption{Acteurs du système PMS et leurs limites}
\label{tab:acteurs}
\end{longtable}

\section{Besoins fonctionnels}
Les besoins fonctionnels sont organisés par module. Ils synthétisent les exigences du document SRS
(FR-001 à FR-053) et les cas d'utilisation (UC-001 à UC-020).

\begin{longtable}{|m{3.5cm}|m{11cm}|}
\hline
\textbf{Module} & \textbf{Principaux besoins fonctionnels} \\
\hline
\endhead
\endfoot
\endlastfoot
\hline
Authentification & Connexion par e-mail / mot de passe, génération de jetons JWT et de
rafraîchissement, changement obligatoire du mot de passe à la première connexion, déconnexion. \\
\hline
Gestion des utilisateurs & Création par l'administrateur, génération automatique d'un mot de passe
temporaire, modification, activation / désactivation, attribution d'un rôle. \\
\hline
Rôles et permissions & Gestion des rôles, des permissions et des associations rôle--permission. \\
\hline
Gestion des projets & Création, modification, changement de statut, définition du budget,
affectation d'un chef de projet, consultation du portefeuille. \\
\hline
Gestion des équipes & Affectation et retrait de développeurs, consultation des disponibilités,
historisation des affectations. \\
\hline
Gestion du TCC & Saisie du coût chargé par ressource et par année, historisation, calcul du coût
journalier. \\
\hline
Plan de charge & Planification mensuelle en jours-homme, totaux planifiés, calcul de l'ETC et des
prévisions. \\
\hline
Charges réelles & Saisie des jours-homme réalisés (projets affectés uniquement), validation,
historisation, import depuis KIMAI. \\
\hline
Facturation & Création des jalons (pourcentage, montant, dates), suivi des statuts, calcul de
l'avancement. \\
\hline
Missions & Création des missions, saisie des frais, calcul automatique du coût total. \\
\hline
KPI et reporting & Calcul des indicateurs financiers, tableaux de bord projet et portefeuille,
génération et export de rapports. \\
\hline
Gouvernance (avenants, risques, livrables, parties prenantes, changements) & Registres de suivi de
la vie du projet, conformément au classeur Excel. \\
\hline
\caption{Besoins fonctionnels par module}
\label{tab:besoins-fonctionnels}
\end{longtable}

\section{Besoins non fonctionnels}
\begin{itemize}
    \item \textbf{Performance}~: réponse en moins de deux secondes pour les opérations courantes
    (NFR-001).
    \item \textbf{Sécurité}~: authentification JWT obligatoire, mots de passe chiffrés (NFR-002).
    \item \textbf{Autorisation}~: contrôle d'accès basé sur les permissions (RBAC dynamique,
    NFR-003).
    \item \textbf{Maintenabilité}~: architecture propre, modulaire et orientée fonctionnalités
    (NFR-006).
    \item \textbf{Évolutivité}~: support de la croissance et intégration future (ERP, RH,
    comptabilité) sans refonte du cœur (NFR-005, BR-065).
    \item \textbf{Auditabilité}~: conservation de l'historique et suppression logique (NFR-008).
    \item \textbf{Ergonomie}~: interface intuitive, responsive et en langue française (NFR-007).
\end{itemize}

\section{Règles de gestion}
Les règles de gestion proviennent du document BRS (65 règles, BR-001 à BR-065), enrichies par les
règles extraites du fichier Excel. Les plus structurantes sont synthétisées ci-dessous.

\begin{longtable}{|m{3.5cm}|m{11cm}|}
\hline
\textbf{Domaine} & \textbf{Règles principales} \\
\hline
\endhead
\endfoot
\endlastfoot
\hline
Utilisateurs & Création réservée à l'administrateur~; pas d'auto-inscription~; un seul rôle par
utilisateur~; e-mail unique~; mot de passe temporaire puis changement à la première connexion. \\
\hline
Projets & Création réservée au directeur~; exactement un directeur et un chef de projet actif par
projet~; statut parmi \{DRAFT, ACTIVE, ON\_HOLD, COMPLETED, CANCELLED\}~; code projet unique~;
budget strictement positif~; date de début antérieure à la date de fin. \\
\hline
Équipes & Un développeur peut appartenir à plusieurs projets~; le retrait conserve l'historique~;
l'affectation enregistre son auteur et sa date~; les développeurs inactifs ne peuvent être
affectés. \\
\hline
TCC & Géré par l'administrateur~; une valeur par ressource et par année~; valeur historique
immuable après validation financière. \\
\hline
Charges & Charges exprimées en jours-homme, jamais négatives~; saisie réservée aux projets
affectés~; historique préservé (aucune suppression). \\
\hline
Facturation & Jalon rattaché à un projet~; statut parmi \{PENDING, INVOICED, PAID\}~; montant
positif~; somme des pourcentages des jalons inférieure ou égale à 100~\%. \\
\hline
Missions & Mission rattachée à un projet et à un employé~; date de début antérieure à la date de
fin~; coût total calculé automatiquement. \\
\hline
KPI et sécurité & KPI financiers visibles uniquement par le directeur et le chef de projet~;
calculs automatiques~; RBAC appliqué par permissions~; mots de passe chiffrés. \\
\hline
\caption{Synthèse des règles de gestion}
\label{tab:regles-gestion}
\end{longtable}

\section{Le modèle d'autorisation dynamique}
\label{sec:rbac}
L'exigence architecturale centrale du projet est de \textbf{ne jamais coupler directement un rôle à
un module}. Un tel couplage (\og{}~ROLE~$\rightarrow$~MODULE~\fg{}) imposerait de modifier le code
source à chaque évolution des responsabilités. Le système adopte donc un modèle dynamique à trois
niveaux~:
\[
\text{ROLE} \;\rightarrow\; \text{PERMISSION} \;\rightarrow\; \text{MODULE}
\]

Les rôles sont des ensembles de permissions. L'autorisation côté serveur, la visibilité côté client
et la \textbf{génération dynamique des menus} reposent toutes sur les permissions, jamais sur le nom
du rôle. Le tableau~\ref{tab:permissions} présente l'attribution par défaut des principales
permissions.

\begin{longtable}{|m{6cm}|m{8.5cm}|}
\hline
\textbf{Permission (extrait)} & \textbf{Rôles porteurs par défaut} \\
\hline
\endhead
\endfoot
\endlastfoot
\hline
\texttt{MANAGE\_USERS}, \texttt{MANAGE\_ROLES}, \texttt{MANAGE\_PERMISSIONS} & Administrateur \\
\hline
\texttt{MANAGE\_TCC}, \texttt{MANAGE\_PARAMETERS} & Administrateur \\
\hline
\texttt{CREATE\_PROJECT}, \texttt{EDIT\_PROJECT}, \texttt{CHANGE\_PROJECT\_STATUS},
\texttt{MANAGE\_PROJECT\_BUDGET}, \texttt{ASSIGN\_PROJECT\_MANAGER} & Directeur \\
\hline
\texttt{ASSIGN\_DEVELOPER}, \texttt{REMOVE\_DEVELOPER} & Directeur \emph{et} Chef de Projet \\
\hline
\texttt{MANAGE\_WORKLOAD\_PLAN}, \texttt{MANAGE\_BILLING}, \texttt{MANAGE\_MISSIONS} & Chef de
Projet \\
\hline
\texttt{VIEW\_FINANCIALS}, \texttt{VIEW\_KPI} & Directeur (tous projets), Chef de Projet (projets
gérés) \\
\hline
\texttt{VIEW\_EXECUTIVE\_DASHBOARD} & Directeur \\
\hline
\texttt{SUBMIT\_ACTUAL\_WORKLOAD} & Développeur \\
\hline
\caption{Extrait de la matrice Rôle--Permission par défaut}
\label{tab:permissions}
\end{longtable}

\subsection{Illustration~: l'affectation des développeurs}
Les documents de spécification (plan de projet, SRS, cas d'utilisation) indiquent que seul le chef
de projet peut affecter des développeurs. Le besoin métier validé exige cependant que le
\emph{directeur} puisse également le faire. Plutôt que de modifier le code, il suffit d'attribuer la
permission \texttt{ASSIGN\_DEVELOPER} au rôle \texttt{DIRECTEUR}~:
\begin{center}
Aujourd'hui~: \texttt{CHEF\_DE\_PROJET} $\rightarrow$ \texttt{ASSIGN\_DEVELOPER} \\
Demain~: on ajoute \texttt{DIRECTEUR} $\rightarrow$ \texttt{ASSIGN\_DEVELOPER}
\end{center}
Ce changement s'opère par une simple mise à jour de configuration en base de données, sans aucune
modification du \emph{backend}, du \emph{frontend} ni de la logique métier. Cet exemple illustre
concrètement la valeur du modèle d'autorisation dynamique.

\section{Le modèle financier et les indicateurs de performance}
Le cœur métier de la plateforme réside dans son modèle financier, fidèle à la réalité extraite du
classeur Excel (onglets \emph{TCC}, \emph{Coût Prévisionnel}, \emph{Coût réel Actualisé},
\emph{Avancement de facturation}, \emph{Missions} et \emph{Glossaire}). Les constantes métier
(taux de frais généraux, jours ouvrés, provisions) sont paramétrables et jamais codées en dur.

\begin{longtable}{|m{4cm}|m{10.5cm}|}
\hline
\textbf{Indicateur} & \textbf{Définition / Calcul} \\
\hline
\endhead
\endfoot
\endlastfoot
\hline
TCC journalier & (coût annuel chargé $\times$ (1 + taux de frais généraux)) / jours ouvrés de
l'année. Dans l'Excel~: facteur $\times 1{,}7$ (frais généraux $\approx 70$~\%) et 22 jours (2024)
ou 20 jours (2025). \\
\hline
Coût Prévisionnel & somme, par mois, des jours-homme planifiés $\times$ TCC. \\
\hline
ETC (reste à faire) & somme des jours-homme restants $\times$ TCC. \\
\hline
Coût Actuel & somme des jours-homme consommés $\times$ TCC. \\
\hline
EAC & Coût Actuel $+$ ETC. \\
\hline
Earned Value (EV~\%) & taux d'avancement réel du projet. \\
\hline
Cumul CA Production & total du contrat $\times$ EV~\%. \\
\hline
FAE / Stock & Cumul CA Production $-$ total facturé. \\
\hline
RAF (jours-homme) & charge planifiée restante $-$ charge consommée. \\
\hline
Dérive (jours-homme) & charges vendues $-$ charges consommées $-$ RAF. \\
\hline
Marge nette vendue~\% & (prix de vente $-$ coût total) / prix de vente. \\
\hline
Marge nette actuelle & Cumul CA Production $-$ Coût Actuel. \\
\hline
Marge nette EAC & Budget total $-$ (Coût Actuel $+$ Coût Prévisionnel), soit Budget $-$ EAC. \\
\hline
Avancement de facturation~\% & cumul facturé / total du contrat. \\
\hline
Coût de mission & perdiem (séjour $\times$ indemnité $\times$ taux de change) $+$ billet $+$ timbre
($+$ transport, séjour). \\
\hline
\caption{Modèle financier et indicateurs de performance (modèle Excel complet)}
\label{tab:kpi}
\end{longtable}

\subsection{KPI par rôle}
Le développeur ne dispose d'aucun indicateur financier~: il ne consulte que ses propres charges. Il
en fournit néanmoins la mesure brute, car sa saisie des jours-homme réels alimente la chaîne
Coût Actuel $\rightarrow$ EAC $\rightarrow$ Marge. Le chef de projet exploite l'ensemble des KPI des
projets qu'il gère afin de piloter leur rentabilité, tandis que le directeur dispose d'une vue
consolidée du portefeuille via un tableau de bord exécutif. L'administrateur, enfin, n'accède à
aucun KPI projet mais administre les intrants (TCC, paramètres) dont ils dépendent.

\section{Processus métier global}
Le déroulement nominal de la plateforme enchaîne~: création des comptes par l'administrateur~;
première connexion et changement de mot de passe~; création d'un projet et affectation d'un chef de
projet par le directeur~; constitution de l'équipe~; planification de la charge~; définition des
jalons de facturation et des missions~; saisie des charges réelles par les développeurs (ou import
KIMAI)~; calcul automatique des KPI~; supervision du portefeuille par le directeur.

\section{Analyse du classeur Excel existant}
La source métier principale est le classeur de revue de projet
(\textit{DEV\_Fiche\_revue\_projet}), composé de \textbf{21 onglets} et de \textbf{2\,088 formules}.
Chaque onglet a été analysé (objet, propriétaire métier, règles, calculs) et rattaché à un module
de la future plateforme. Le tableau~\ref{tab:excel-onglets} en présente la synthèse.

\begin{longtable}{|m{4cm}|m{4cm}|m{6cm}|}
\hline
\textbf{Onglet} & \textbf{Propriétaire métier} & \textbf{Module PMS} \\
\hline
\endhead
\endfoot
\endlastfoot
\hline
Fiche Doc & Qualité / Direction & Gestion de projet (métadonnées) \\ \hline
Fiche identification & Directeur & Gestion de projet \\ \hline
TCC & Administrateur & Gestion du TCC \\ \hline
Missions & Chef de projet & Gestion des missions \\ \hline
Avancement\_Facturation & Chef de projet & Facturation \\ \hline
Recap Facturation (2) & Chef de projet & Facturation \\ \hline
Coût Prévisionnel & Chef de projet & Plan de charge + KPI \\ \hline
Coût réel Actualisé & Développeur / Chef de projet & Charges réelles + KPI \\ \hline
Recap Facturation & Chef de projet & Facturation \\ \hline
S. Contractuelle & Directeur / Chef de projet & Gestion de projet + Facturation \\ \hline
Situation actuelle & Chef de projet & KPI / Livrables \\ \hline
Parties Prenantes & Chef de projet & Parties prenantes + Équipe \\ \hline
Livrables projet & Chef de projet & Livrables \\ \hline
Base des Risques & Chef de projet & Registre des risques \\ \hline
Liste des actions & Chef de projet & Actions / Risques \\ \hline
Registre Changement & Chef de projet & Registre des changements \\ \hline
Avenants & Directeur / Chef de projet & Avenants \\ \hline
Statistiques & Système (calcul) & Moteur de KPI \\ \hline
Dashboard & Système / Directeur, CP & Tableaux de bord \\ \hline
Glossaire & Référence & KPI (définitions) \\ \hline
Paramètres & Administrateur & Configuration / Échelles de risque \\ \hline
\caption{Analyse des 21 onglets du classeur Excel et rattachement aux modules PMS}
\label{tab:excel-onglets}
\end{longtable}

\subsection{Couverture des formules}
Les 2\,088 formules du classeur reposent sur un vocabulaire \emph{fini}, ce qui permet de garantir
que chacune est rattachée à un calcul futur du système. Les catégories ci-dessous sont exhaustives.

\begin{longtable}{|m{6.5cm}|m{2cm}|m{5.5cm}|}
\hline
\textbf{Catégorie de formule} & \textbf{Nombre} & \textbf{Calcul PMS associé} \\
\hline
\endhead
\endfoot
\endlastfoot
\hline
Fonctions de date (DATE, YEAR, MONTH, DAY, EDATE, TODAY) & 460 & Calendrier projet/sprint, sélection du mois courant, durées \\ \hline
Agrégation SUM & 92 & Totaux mensuels et cumuls (coûts, facturation, charges) \\ \hline
HLOOKUP (sélection du mois) & 31 & Instantané KPI du mois sélectionné \\ \hline
IF / ISBLANK / IFERROR & 36 & Statuts conditionnels, valeurs sûres \\ \hline
AVERAGE / MAX / COUNTA & 7 & Moyennes, maximum, comptage de livrables \\ \hline
Arithmétique de coûts (JH $\times$ TCC, sommes, écarts) & 1\,421 & Coût prévisionnel, coût actuel, ETC, EAC, coût mission \\ \hline
Pourcentages & 13 & Marges, avancement de facturation, PPR, Delivery \% \\ \hline
Références inter-onglets & 1\,426 & Jointures relationnelles (charges $\leftrightarrow$ TCC $\leftrightarrow$ projet) \\ \hline
Références cassées (\#REF!, \#VALUE!) & 169 & Non reproduites : recalcul propre à partir de données normalisées \\ \hline
\caption{Cartographie des formules Excel vers les calculs du futur système}
\label{tab:formules}
\end{longtable}

\noindent\textbf{Résultat de la vérification :} chacune des 2\,088 formules relève d'une catégorie
de logique métier vivante (rattachée à un calcul PMS) ou d'un artefact cassé que la plateforme
remplace par un recalcul propre. \emph{Aucune formule n'est laissée sans correspondance.}

\section*{Conclusion}
Ce chapitre a spécifié les acteurs, les besoins, les règles de gestion, le modèle d'autorisation
dynamique, le modèle financier et l'analyse complète du classeur Excel existant. Ces éléments
définissent sans ambiguïté le \emph{quoi} du système, strictement sur le plan métier. Les chapitres
suivants aborderont le \emph{comment}~: l'architecture technique, la conception de la base de
données, la modélisation UML puis la réalisation des modules.
